\chapter{Enveloping algebras and planar-sequence obstructions}
\label{ch:enveloping-audit}

\section{The bar of an enveloping algebra}

\begin{theorem}[Bar--Chevalley comparison]
\label{thm:bar-chevalley}
\StatusProved{} Let $\mathfrak g$ be a Lie algebra object in $\FactD(X)$
for the pointwise tensor, $K$-flat, positively graded and complete, with
every total weight admitting finitely many decompositions, and suppose
internal PBW holds for $U(\mathfrak g)$. There is a natural
quasi-isomorphism of conilpotent coalgebras in $\FactD(X)$
\[
  \Barperp_X\,U(\mathfrak g)\ \xrightarrow{\ \sim\ }\ \CE_*(\mathfrak g),
\]
compatible with the factorisation coefficient and the PBW filtration.
\end{theorem}

\begin{proof}
Filter $U(\mathfrak g)$ by PBW degree and both complexes by total word
length. The associated graded enveloping algebra is $\Sym(\mathfrak g)$,
and the bar of a symmetric algebra is the suspended symmetric coalgebra on
$\mathfrak g$, which is the associated graded of Chevalley chains. The
comparison is an isomorphism on the first page; finite decomposition gives
a bounded exhaustive spectral sequence in each weight and completeness
gives the quasi-isomorphism. Every step uses only internal tensor products
and the bracket, so the map stays in $\FactD(X)$.
\end{proof}

Theorem~\ref{thm:bar-chevalley} is the only route by which a Lie-theoretic
calculation computes an ordered bar, and it fixes what may be tabulated:
the Chevalley--Eilenberg homology of the Lie algebra, in the grading the
Lie algebra carries.

\section{Simple Lie algebras}

\begin{computation}[Finite-dimensional simple case]
\label{comp:simple}
\StatusComputed{} For a simple Lie algebra $\mathfrak g$ of rank $r$ with
exponents $m_1,\dots,m_r$, $H^*(\mathfrak g;k)$ is the exterior algebra on
primitive classes of degrees $2m_i+1$. Hence
\[
  \Barperp U(\mathfrak{sl}_2):\ 1+t^3, \qquad
  \Barperp U(\mathfrak{sl}_3):\ (1+t^3)(1+t^5) = 1+t^3+t^5+t^8,
\]
\[
  \Barperp U(\mathfrak g_2):\ (1+t^3)(1+t^{11}) = 1+t^3+t^{11}+t^{14}.
\]
The first two were recomputed here by exact rational elimination on
$\Lambda^\bullet\mathfrak g^*$ built from matrix units, with $d^2=0$
asserted on every cochain space rather than assumed; the results are
$\dim H^n(\mathfrak{sl}_2) = 1,0,0,1$ and
$\dim H^n(\mathfrak{sl}_3) = 1,0,0,1,0,1,0,0,1$. The chain dimensions are
$\binom{\dim\mathfrak g}{n}$, so for $\mathfrak{sl}_2$ they are
$1,3,3,1$.
\end{computation}

\section{The Virasoro family}

The transverse presentation of the Virasoro family is the enveloping
algebra of the positive Witt algebra
\[
  L_1 \;=\; \bigoplus_{n\ge1}\mathbb C\,e_n,
  \qquad [e_i,e_j]=(j-i)\,e_{i+j},
\]
to which Theorem~\ref{thm:bar-chevalley} applies: $L_1$ is positively
graded with one-dimensional weight spaces, so every weight has finitely
many decompositions.

\begin{theorem}[Pentagonal rigidity]
\label{thm:pentagonal-rigidity}
\StatusComputed{} The bar homology of $U(L_1)$ satisfies
\[
  \dim H^n\bigl(\Barperp U(L_1)\bigr) = 2
  \qquad\text{for every } n\ge1,
\]
one dimension in each of the weights $(3n^2-n)/2$ and $(3n^2+n)/2$. The
graded Euler character is
\begin{equation}
\label{eq:witt-euler}
  \chi_q\bigl(\Barperp U(L_1)\bigr) = \prod_{n\ge1}\bigl(1-q^n\bigr),
\end{equation}
and it is accounted for degree by degree with no residual freedom.
\end{theorem}

\begin{proof}
By Theorem~\ref{thm:bar-chevalley} the bar of $U(L_1)$ is
$\CE_*(L_1)$, whose cohomology is Goncharova's \cite{Gon73}: two
dimensional in each degree $n\ge1$, in the weights $(3n^2\mp n)/2$.
Equation~\eqref{eq:witt-euler} is Theorem~\ref{thm:bar-denominator}
applied to $L_1$, where every root space is one dimensional and even, so
$\sdim (L_1)_n = 1$. Euler's pentagonal number theorem expands the
product as $\sum_{m\in\mathbb Z}(-1)^mq^{m(3m-1)/2}$, whose support is
exactly $\{(3n^2\pm n)/2\}$ with coefficient $(-1)^n$ at each of
$(3n^2-n)/2$ and $(3n^2+n)/2$. Comparing with $\sum_n(-1)^n\dim H^n_w$
shows each weight of the Euler character is saturated by a single class in
the degree Goncharova predicts, so no further classes and no cancellation
are available.
\end{proof}

The cohomology was computed for $n=1,2,3,4$, occupying the weights
$(1,2)$, $(5,7)$, $(12,15)$, $(22,26)$, one dimension in each; the nonzero
coefficients of $\prod_{n\ge1}(1-q^n)$ occupy $0,1,2,5,7,12,15,22,26$ with
signs $+,-,-,+,+,-,-,+,+$.

\section{The proposed planar sequences}

Two formulas were tabulated as bar cohomology dimensions:
\[
  \dim H^n = M(n+1)-M(n) = 1,2,5,12,30,76,196,512,1353,3610,\dots
\]
for the Virasoro family, with $M$ the Motzkin numbers, and
\[
  \dim H^n = R(n+3) = 3,6,15,36,91,232,603,1585,\dots
\]
for the $\mathfrak{sl}_2$ family, with $R$ the Riordan numbers. The
coincidence of the discriminant $\sqrt{1-2z-3z^2}$ in the two generating
functions was attributed to Drinfeld--Sokolov reduction.

\begin{proposition}[The two sequences are one sequence]
\label{prop:one-sequence}
\StatusComputed{} For all $n\ge0$,
\[
  M(n) = R(n)+R(n+1),
  \qquad
  M(n+1)-M(n) = R(n+2)-R(n).
\]
The Virasoro sequence is therefore $R(n+2)-R(n)$ and the
$\mathfrak{sl}_2$ sequence is $R(n+3)$: both are elementary transforms of
one sequence, and their generating functions are rational in the single
algebraic function $\sqrt{1-2z-3z^2}$ for that reason alone. No
Drinfeld--Sokolov input is used or needed.
\end{proposition}

\begin{proof}
Riordan numbers count Motzkin paths with no level step at height zero.
Splitting a Motzkin path of length $n$ according to whether its first step
is level gives $M(n)=R(n)+R(n+1)$; subtracting consecutive instances gives
the second identity. Both were checked termwise for $n<24$.
\end{proof}

\begin{theorem}[Non-identification with ordered bar homology]
\label{thm:exact-audit}
\StatusComputed{} Neither sequence is the ordered bar homology of the
algebra it names, and neither is a bar chain count.
\begin{enumerate}
\item For $\mathfrak{sl}_2$: the bar homology is $1,0,0,1$ by
  Computation~\ref{comp:simple}, and for $\dim\bar A = 3$ the length-$r$
  ordered bar term has dimension $3^r$, giving $3,9,27,81,243$. The
  sequence $3,6,15,36,91$ is neither.
\item For the Virasoro family: the bar homology is the constant $2$ by
  Theorem~\ref{thm:pentagonal-rigidity}. On degrees $n=1,2,3$ the asserted
  values are $1,2,5$; they agree with the truth at $n=2$ and nowhere else.
\item Independently of any presentation, a Virasoro chiral coefficient
  does not determine the transverse bar at all: a fixed coefficient
  $\mathcal F_{\mathrm{Vir},c}$ may be combined with free, square-zero,
  enveloping, or quantum-plane transverse algebras whose bars differ,
  while the coefficient and the central charge are unchanged.
\end{enumerate}
\end{theorem}

\begin{proof}
(i) and (ii) are the cited computations. For (iii), form $B\otimes^!
\mathcal F_{\mathrm{Vir},c}$ for the four transverse algebras listed; the
coefficient is fixed and $\Barperp$ varies with $B$.
\end{proof}

\begin{remark}[The type error beneath the numerical one]
\label{rem:type-error}
The bar homology of $U(L_1)$ is bigraded by cohomological degree and
weight, and is finite in each degree only because
Theorem~\ref{thm:pentagonal-rigidity} concentrates it in two weights. A
formula assigning one integer to each degree, with no weight truncation
named, does not specify an invariant of a graded object with
infinite-dimensional graded pieces. The sequences above are counts of
planar objects; passage from such a count to a homology dimension requires
a differential, and no quasi-isomorphism to a Motzkin or Riordan complex
has been constructed.
\end{remark}

\begin{remark}[Three distinct chain complexes]
\label{rem:three-models}
The interval bar $\Barperp(A)$, whose
differential uses transverse multiplication; chiral Chevalley chains,
whose differential uses a diagonal-supported bracket; and a
screen-enhanced bicomplex $B^{p,q}_{\mathrm{scr}}(A) = (s\bar
A)^{\otimes^! p}\otimes C^q_{\mathrm{scr}}$ with total differential
$b_A+d_A+d_{\mathrm{scr}}+\delta_{\mathrm{int}}$ are distinct complexes.
Only the third can carry both bar words and
configuration-space cohomology, and only after
$\delta_{\mathrm{int}}$ is supplied and shown to square to zero. Filtering
by screen degree gives $E_1^{p,q} = H^q(C_{\mathrm{scr}})\otimes
B^\perp_p(A)$; a product of dimensions on an $E_1$ page is not the
dimension of total homology.
\end{remark}
